\documentclass[11pt]{article}
\usepackage[margin=1.1in]{geometry}
\usepackage{amsmath,amssymb,amsthm}
\usepackage{booktabs}
\usepackage{graphicx}
\usepackage{enumitem}
\usepackage[colorlinks=true,linkcolor=blue,citecolor=blue,urlcolor=blue]{hyperref}
\usepackage{xcolor}
\usepackage{listings}

% ---- Notation ----
\newcommand{\X}{\chi}
\newcommand{\sep}{\ensuremath{\mathtt{0x1E}}}
\newcommand{\rparam}{r}
\newcommand{\Order}{\mathcal{O}}
\newcommand{\Bits}{$\mathrm{bits}$}
\newcommand{\TODO}[1]{\textbf{\textcolor{red}{[TODO: #1]}}}
\newcommand{\hprcchi}{\TODO{466 headline $\X$ lands with the running build; expected $\approx 2.25\times 10^9$}}

% Theorem environments
\newtheorem{theorem}{Theorem}[section]
\newtheorem{lemma}[theorem]{Lemma}
\newtheorem{corollary}[theorem]{Corollary}
\newtheorem{definition}[theorem]{Definition}
\newtheorem{remark}[theorem]{Remark}
\newtheorem{conjecture}[theorem]{Conjecture}

% Alternate title (kept per lane record): "Indexing Novelty, Not Bulk:
% The Suffixient Array at Pangenome Scale, with a First Space Lower Bound
% for Pattern Matching"
\title{$\chi$ Marks the Fork:\\
Suffixient Arrays and a Space Lower Bound for Locating Patterns}

\author{%
  E.~Gaffney\thanks{Correspondence: the sxgc ledger, \texttt{github.com/pangenome/sxgc}.}%
}
\date{September 2026}

\begin{document}
\maketitle

\begin{abstract}
\textbf{Motivation.} Almost every byte of a corpus is \emph{forced}:
given its left context, only one continuation occurs. The freedom lives in a
sparse skeleton of \emph{forks}---positions where the future is genuinely
open---and at pangenome scale that skeleton is tiny: measured across this
paper's corpora, the minimum suffixient set $S$ (a set of positions witnessing
every right-extension context of the text) has $\X/n$ at $2.6\%$ (yeast),
$5.4\%$ (30\,Gbp human), and $0.16\%$ (466-haplotype HPRC, $1.44$\,Tbp).
The suffixient array \cite{measure-paper} measures this skeleton: $\X(T)$
upper bounds the string attractor measure $\gamma$ and is at most twice the
run count of the BWT of the reversed text. But $\X$ was a \emph{one-sided}
story: it said how small an index \emph{can} be, never how small any index
\emph{must} be---no space lower bound proportional to any repetitiveness
measure existed for any pattern-matching operation class (verified survey,
\S\ref{sec:related}). Three questions were open when this work began: \emph{can
$\X$ be constructed} without reading the text more than once and without any
operation superlinear in $n$; \emph{is $\X$ necessary}; and \emph{does the
theory build an index} at terabase scale.
\textbf{Results.} We give a machine-checked (Lean 4) development in which
$\X(T)$ equals the number of inclusion-maximal coverage classes of the text's
requirements, a one-pass scan over the text's suffix-array stream emits a
minimum suffixient set (the construction capstone, proven subject to named,
explicitly-tracked stream hypotheses: the no-duplicates hypothesis is discharged outright, and the remaining hypotheses reduce to a single named statement about run-edge domination, one of whose two factors is proven), the
parse-space LCE machinery the scan consumes is proven deterministic and
explicitly bounded, and the head-sample extraction used at scale is proven as
an inverse permutation on BWT rows. We then make $\X$ two-sided. \emph{Bridge}:
any system---index, language model, anything---that correctly answers
\emph{continuation} queries (``what can follow $w$'', answered with corpus
positions) must consult at least $\X(T)$ \emph{distinct} positions
(kernel-checked; this form strictly subsumes the locate-one bridge, since
locate-one answers are correct continuation answers). \emph{Floor}: on a
constructed witness-perturbation family, every fixed-decoder index answering
locate-one satisfies $s+1 \ge \X\log_2(n/\X)/3$---to our knowledge the first
space lower bound for any pattern-matching operation class parametrized by a
repetitiveness measure, stated with its honest scope (family-specific,
fixed-decoder). We describe a system that constructs the index as one command
from a compressed pangenome archive, never materializes the text on any disk,
and publishes only after an in-flight audit of sampled witnesses against the
archive. At yeast scale the index is born from scratch in 26 minutes; at
30\,Gbp human scale, $\X = 1{,}627{,}067{,}257$ with peak 134\,GB; at the full
466-haplotype HPRC collection (1.44\,Tbp) the headline number is \hprcchi.
Every claim in this paper is traceable to a machine-checked proof or a logged
measurement; we state the open hypotheses and unproven attributions exactly
where they remain.
\textbf{Availability.} Proofs, logs, and the tool at
\texttt{github.com/pangenome/sxgc}.
\end{abstract}

\section{Introduction}\label{sec:intro}

\paragraph{Most bytes are forced.} Measure any real corpus: for the vast
majority of positions, the byte that follows is the only one that ever does.
Across this paper's corpora the \emph{minimum suffixient set}---the smallest
position set witnessing every place where the future is genuinely open---has
$\X/n$ at $2.6\%$ for yeast, $5.4\%$ for a 30\,Gbp human pangenome, and
$0.16\%$ for the 466-haplotype HPRC collection
($1.44$\,Tbp; \S\ref{sec:results}). The farther a collection duplicates
itself, the smaller the skeleton: a pangenome is $99.84\%$ inevitability.
Whatever an exact-search index---or any reader of the corpus---must retain is
therefore a function of this sparse forking structure, not of the bulk.

\paragraph{The thesis.} Suppose a corpus is mostly redundant---as every
versioned corpus is. Then whatever an exact-search index must store is a
function of the corpus's \emph{novelty}, not its bulk. This paper is about
making that sentence a theorem-backed engineering statement: a measure
$\X(T)$ that quantifies novelty, a construction that builds an index of size
$O(\X\log n)$ at a cost that follows $\X$ rather than $n$, a lower-bound
program that has produced its first theorem saying that position-answering
indexes \emph{cannot} be asymptotically smaller, a bridge theorem extending
that necessity from indexes to \emph{any} continuation-answering system
(\S\ref{sec:contbridge}), and a system that does all of it from a 3\,GB
compressed archive of 1.44 terabases without ever writing the text to disk.

\paragraph{The object.} A \emph{suffixient set} for a text $T$ of length $n$
is a set $S\subseteq[1..n]$ of positions such that for every
\emph{requirement}---a pair $(w,c)$ where $w$ is right-maximal and $wc$ occurs
in $T$---some $x\in S$ has $wc$ as a suffix of the prefix of $T$ ending at
$x$. The smallest such set has cardinality $\X(T)$, introduced by
Fujimaru, Navarro, Romana and Urbina~\cite{measure-paper}, who proved
$\gamma \le \X \le 2\overline r$ (where $\overline r$ is the run count of the
BWT of the reversed text), showed that a suffixient set doubles as a string
attractor and hence supports $O(\X\log(n/\X))$-space text storage, and gave
calculation methods. What was \emph{not} known: whether $\X$ is the minimum
\emph{size of a certain kind of index} (they prove minimality among suffixient
sets, not among indexes), whether $\X$ is computable at scale under the
one-pass discipline realistic constructions obey, and whether any space lower
bound proportional to \emph{any} repetitiveness measure holds for any
pattern-matching class (a verified literature survey we summarize in
\S\ref{sec:related} found none: every published $r$- or $\X$-index is an upper
bound whose ``optimality'' is time-only).

\paragraph{The program.} We run a three-legged program in which no leg trusts
the others: a \emph{formal} leg (Lean 4) that proves the definitions,
characterization, construction, bounds, and now a first lower bound; an
\emph{empirical} leg in which brute-force differentials on small texts are
ground truth and every large-scale claim is gated by byte-exact comparison
against an independently built oracle; and a \emph{systems} leg that
implements the construction as a one-command pipeline and records, per
stage, wall time and peak memory from the machine. The discipline that binds
them is stated plainly: statements are locked and measurement-owned; proofs
are kernel-checked with a fixed axiom base; every numeric claim in this paper
cites the artifact that produced it (Appendix~\ref{app:facts}).

\paragraph{Contributions.}
\begin{enumerate}[leftmargin=*]
\item \textbf{Characterization} (\S\ref{sec:formal}): $\X(T)$ equals the number
of inclusion-maximal coverage classes of $T$'s requirements
(Theorem~\ref{thm:maxclasses}, machine-checked), which turns the brute-force
minimum into a structural count and is the engine behind the scan.
\item \textbf{Construction} (\S\ref{sec:construction}): a one-pass scan over
the text's suffix-array stream (BWT char, LCP, SA) emits a minimum suffixient
set; the scan is proven nodup for all well-formed bounded streams
(Theorem~\ref{thm:nodup}); the parse-space LCE machinery the scan consumes at
scale is proven deterministic (two-sided equality with ground truth) and
explicitly bounded (Theorems~\ref{thm:bounds}--\ref{thm:totalwork}). The
end-to-end statement---the parse-built witness list has exactly $\X(T)$
elements---is proven \emph{conditional on named stream hypotheses}
(Theorem~\ref{thm:capstone}), of which all but one are discharged; the
remaining one is reduced to a single statement about run-edge domination
(``RunEdgeDominate'', \S\ref{sec:status}) with its semantic half proven.
\item \textbf{Position extraction} (\S\ref{sec:phi}): the head samples used
to resolve scan positions at scale are the values of an inverse permutation
$\varphi^{-1}$ on BWT rows (Theorem~\ref{thm:headfromtail}), which replaces
the classical LF-step walk---measured at 2{,}156 steps per head on our
corpora---by a lookup.
\item \textbf{Cyclic frames} (\S\ref{sec:seam}): the padded cyclic BWT the
upstream front end produces and the canonical cyclic order of the corpus text
agree \emph{outside seam classes}; class repair is the identity there and a
bounded in-class reordering. The identity half is machine-checked
(Theorem~\ref{thm:seam}); the repair is implemented with an explicit refusal
policy for adversarial corpora.
\item \textbf{A first space lower bound} (\S\ref{sec:floor}): every
fixed-decoder index answering locate-one queries over the
witness-perturbation family of texts satisfies
\[
s + 1 \;\ge\; \X(T)\cdot \log_2\!\big(n/\X(T)\big)\,/\,3 ,
\]
machine-checked (Theorem~\ref{thm:floor}). To our knowledge this is the first
space lower bound proportional to a repetitiveness measure for any
pattern-matching operation class (survey in \S\ref{sec:related}).
\item \textbf{The continuation bridge} (\S\ref{sec:contbridge}): the same
necessity, extended from indexes to \emph{any} system that answers
continuation queries. Any oracle that, for every requirement $(w,c)$ of $T$,
answers with a genuine corpus position where $wc$ is covered must \emph{consult}
at least $\X(T)$ \emph{distinct} positions over the course of answering
(kernel-checked, Theorem~\ref{thm:contdistinct}); the demand is strictly
weaker than locate-one, so the new bridge subsumes the old
(Theorem~\ref{thm:locimpcont}). This is the form the language-model
framing of \S\ref{sec:lm} consumes.
\item \textbf{System and measurements} (\S\ref{sec:system},
\S\ref{sec:results}): a one-command pipeline that parses directly from the
compressed archive (no text file is ever written), repairs the frame, extracts
samples by $\varphi^{-1}$, computes aggregates in parse space, sweeps the
witnesses, audits 10{,}000 of them in flight against the archive, and
publishes a single container file holding the run-length BWT, both SA-sample
sets, names, and $\delta$-compressed $\X$ ($\approx 13\%$ of raw). Measured
end-to-end at yeast (26 minutes, 8.8\,GB peak) and at 30\,Gbp human scale
($\X=1{,}627{,}067{,}257$, 134\,GB peak); the 466-haplotype HPRC run
is \hprcchi.
\end{enumerate}

\section{The Formal Model}\label{sec:formal}

All statements in this section are stated as they appear in the Lean sources
(cited in Appendix~\ref{app:facts}); the axiom base is exactly
\{\texttt{propext}, \texttt{Classical.choice}, \texttt{Quot.sound}\}. The
text is a list of naturals $T$; positions are $1$-based.

\begin{definition}[Requirements; \texttt{Sxgc.lean}]
A pair $(w,c)$ is a \emph{requirement} of $T$ if $w$ is right-maximal in $T$
and $wc$ occurs in $T$, where right-maximal means $w$ occurs and either $w$
is a suffix of $T$ or $w$ has $\ge 2$ distinct right extensions. Position $x$
\emph{covers} $(w,c)$ if $wc$ is a suffix of the prefix of $T$ of length $x$.
$S$ is \emph{suffixient} if every requirement is covered by some $x\in S$.
$\X(T)$ is the minimum cardinality of a suffixient set.
\end{definition}

\begin{definition}[Coverage classes; \texttt{Sxgc.lean}]
For a position $x$, its \emph{coverage set} is the set of requirements it
covers. \emph{Scope comparison} $x \preceq y$ says $y$'s coverage set
contains $x$'s restricted to a common domain (the Lean \texttt{ScopeLe}).
Classes are inclusion-maximal coverage sets; \texttt{maxClassCount} counts
the distinct maxima.
\end{definition}

\begin{theorem}[Characterization; \texttt{chi\_eq\_maxClasses}]\label{thm:maxclasses}
For positive $T$, $\X(T) = \mathrm{maxClassCount}(T)$.
\end{theorem}

The proof's two halves are informative: ($\le$) the earliest representative of
each maximal class is suffixient; ($\ge$) each maximal class has a
\emph{private} requirement---one no other class's position can cover---which
forces one position per class into any suffixient set. The private-requirement
property is exactly the algebraic seed the lower bound of
\S\ref{sec:floor} harvests.

\begin{theorem}[Minimality transfer; \texttt{chi\_le\_of\_suffixient}]
\label{thm:chiLe}
If $S\subseteq[1..n]$ is suffixient for $T$ then $\X(T)\le|S|$.
\end{theorem}

\paragraph{The one-pass scan.} The reference construction is a single left-to-right
pass over the \emph{triple stream} of $T$: for each suffix-array row $k$ of
$T^{R}0$ (the reverse of $T$ with a terminal 0), the triple
$(\mathrm{BWT}[k],\ \mathrm{LCP}[k],\ \mathrm{SA}[k])$. The scan maintains an
\emph{armed candidate} (a tentative witness) and emits it at run boundaries;
the emitted positions form the suffixient set. The Lean \texttt{scan} is a
functional mirror of the deployed C++ scanner, byte-verified against it on
all battery corpora.

\begin{theorem}[Scan emits run edges; \texttt{scan\_emits\_run\_edges}]
Every emitted position is the position of a run-edge row of the triple
stream.
\end{theorem}

\begin{theorem}[No duplicates, all bounded streams; \texttt{scanAux\_nodup}]
\label{thm:nodup}
For \emph{any} triple stream $ts$ (not only those arising from a text) with
distinct SA values, all SA values $<N$, and all LCP values $\le\mathrm{MAXINT}$:
the scan's emitted list is duplicate-free.
\end{theorem}

This is the statement-locked pillar \textbf{O2}, now discharged: the proof
threads ten invariants and uses boundedness exactly once, to block re-arming
a just-emitted position at run-length-1 boundaries---the precise point where
unbounded streams duplicate. $1.47$ million bounded test streams showed zero
violations before the proof; the proof replaces the warrant.

\section{The Construction Capstone and Its Ledger}\label{sec:construction}

The construction at scale replaces the text with a prefix-free parse
$(\mathrm{parse},\mathrm{dict})$ and reconstructs the triple stream's LCP
column by two-level fingerprint LCE over the parse. The capstone ties it
together:

\begin{theorem}[Capstone, conditional; \texttt{parseChi\_eq}]\label{thm:capstone}
Let $P$ be a well-formed parse of $T^{R}0$ and let
$S=\mathrm{scan}(\text{parse-triples of }T,P)$. If
(i) $S$ is duplicate-free,
(ii) every $x\in S$ lands in a maximal coverage class (\texttt{IsMax}),
(iii) $S$ is pairwise scope-unique on those classes,
(iv) $S$ is suffixient,
then $|S| = \X(T)$.
\end{theorem}

The four stream hypotheses are named and tracked; their status as of this
writing:

\begin{itemize}[leftmargin=*]
\item[(i)] \emph{Discharged}: Theorem~\ref{thm:nodup} applies to the parse
stream via \texttt{triplesOf\_sa\_nodup}.
\item[(ii),(iii)] \emph{Reduced}: hypothesis (ii) is equivalent to
\texttt{O1\_maxHit}---every maximal class contains an emitted position---and
we prove the factorization
$\mathtt{O1\_maxHit} \iff \mathrm{RunEdgeHit} \wedge \mathrm{RunEdgeDominate}$,
where \texttt{RunEdgeHit} (every maximal class contains a run-edge position)
is \emph{proven outright} (\texttt{SxgcRunEdge.lean}), and
\texttt{RunEdgeDominate} (every run-edge position is dominated by an emitted
one) remains open. Covering (iv) follows from \texttt{maxHit} by
\texttt{covering\_given\_stream\_of\_maxHit} (proven).
\end{itemize}

\textbf{The honest ledger.} \texttt{RunEdgeDominate} is blocked on the exact
emission rule of the scan's LCP-window machinery; its statement-lock
differentials report zero counterexamples over exhaustive
$\{1,2\}^{\le 8}$ and $\{1,2,3\}^{\le 7}$ texts and 3{,}000 random 4-letter
texts, and it is the single remaining open hypothesis of the capstone.

\begin{theorem}[LCE determinism; \texttt{VerifiedLCE\_determinism}]
The parse-space LCE oracle equals ground truth: for a well-formed parse,
$\mathrm{textLCEvia}(P,T,i,j) = \mathrm{lcp}(T[i..], T[j..])$.
\end{theorem}

\begin{theorem}[Explicit bounds; \texttt{SxgcBounds}]\label{thm:bounds}
Each text LCE costs at most
$8\,(1+\tau+\ell)$ elementary steps, where $\tau$ is the fingerprint
sampling period and $\ell$ the number of verified symbols.
\end{theorem}

\begin{theorem}[Total work; \texttt{totalWork\_bound}]\label{thm:totalwork}
For any query batch,
\[
\mathrm{totalWork} \;\le\; 16\Big(\,\mathrm{count} + \mathrm{dictSize}
+ |qs|\cdot\tau + \textstyle\sum_q \ell_q\Big).
\]
For the construction's own queries this yields
$16\,(\mathrm{count}+\mathrm{dict}+n\,\tau+\sum\ell_q)$
(\texttt{parseTriplesLCEWork\_bound}).
\end{theorem}

We deliberately make no claim that the verification lengths $\ell_q$ are
$O(\tau)$; they are measured (15.33 phrases per seed on average, 22{,}262
maximum, on 30\,Gbp human data; \S\ref{sec:results}), and the sums remain
explicit in the theorem.

\section{Position Extraction by Inverse Permutation}\label{sec:phi}

At scale the scan needs, for each run head, a text position; the classical
route is an LF walk from the nearest sampled tail---measured at 2{,}156 steps
per head on our corpora, an $\Omega(n)$ aggregate cost the construction cannot
pay. The alternative is algebraic. Let $\varphi$ map each row to its
predecessor in suffix order (the r-index's $\varphi$). Then:

\begin{theorem}[\texttt{headFromTail}, \texttt{SxgcPhi.lean}]\label{thm:headfromtail}
For a row $a>0$ of the reverse text's index, the head position of row $a$
equals $\varphi^{-1}$ applied to the tail position of row $a-1$:
$\mathrm{rowPos}(a) = \varphi^{-1}(\mathrm{rowPos}(a-1))$.
\end{theorem}

With $\varphi^{-1}$ realized as a lookup over the sorted images of the sample
intervals, every head costs one binary search: on 30\,Gbp human data, all
$1.86\times10^9$ head positions were extracted in 751\,s at 64.8\,GB, versus
a projected 38-hour walk---a $165\times$ replacement, with the extracted
heads byte-identical to the walk baseline on the full corpus.

\section{Cyclic Frames and Seam Classes}\label{sec:seam}

A compressed archive does not define a single text; it defines a family
(the contigs) plus a convention. We adopt the \emph{canonical cyclic frame}:
the corpus is the byte string $T = s_1 \sep s_2 \sep \dots \sep s_k \sep$
($\sep$ a reserved separator byte, $0\le\sep$ below all content bytes), and
the index is its suffix structure. The upstream front end, however, produces
the BWT of the \emph{padded} string $T' = T\$^{w}$ under a cyclic
convention. The two frames agree almost everywhere:

\begin{definition}[Seam position; \texttt{SxgcSeam.lean}]
Row $i$ is a \emph{seam} position if its finite suffix is prefix-comparable
with some other row's suffix.
\end{definition}

\begin{theorem}[Identity outside classes]\label{thm:seam}
For two non-seam rows, finite suffix order and cyclic rotation order agree.
Moreover a non-seam row's order against \emph{every} row is stable, so
in-class reordering can never displace a row across a non-seam row.
\end{theorem}

The repair therefore: strip the padding block (whose structure the front
end certifies), re-sort each seam class by cyclic comparison, splice, and
coalesce. The repair is $r$-space and \emph{fails loudly} rather than pay
$\Theta(n)$ on adversarial (fully periodic) corpora where the class structure
degenerates---our policy is refusal, not a forbidden fallback. Measured:
7 seam classes and 13{,}503 candidate rows at yeast235 ($3.34$\,Gbp,
9{,}901 contigs); 7 classes and 1{,}161 rows at 30\,Gbp human scale. The
in-class reordering lemma itself remains statement-locked with byte-level
warrant (repair output equals a brute-force cyclic oracle on the full
corpora).

\section{The First Space Lower Bound}\label{sec:floor}

No space lower bound proportional to any repetitiveness measure---not $r$,
not $z$, not $\delta$, not $\X$---was previously known for any
pattern-matching operation class (\S\ref{sec:related}). The bridge that makes
one possible for $\X$:

\begin{theorem}[Bridge; \texttt{emitted\_suffixient},
\texttt{chi\_le\_of\_oracle}]\label{thm:bridge}
Let $f$ be any oracle that answers locate-one correctly (it returns, for each
occurring word $wc$ with $(w,c)$ a requirement, a position $x$ with $wc$ a
suffix of $T$'s prefix $x$). Then the set of positions $f$ emits on
$T$'s requirement words is suffixient, and consequently
$\X(T) \le |\mathrm{emitted}(f,T)|$.
\end{theorem}

That is: \emph{any} index, however encoded, that answers ``where can this
continue'' emits at least $\X$ positions. What the bridge does not give is
that those positions must be \emph{stored}: a decoder could emit computed
values. The missing ingredient is a family whose members' correct answers are
pairwise incompatible, so no single decoder can serve two members; then a
counting argument forces one distinct index per member.

\subsection{The Continuation Bridge}\label{sec:contbridge}

The locate-one demand---answer \emph{every} occurring word with an
occurrence---is stronger than the position-answering world of language models
needs. A \emph{continuation oracle} need only answer the requirements
themselves: for each fork $(w,c)$, a genuine corpus position where the
continued word $w{+}{+}[c]$ is covered. The Lean development
(\texttt{LM.lean}; axiom base $\{\texttt{propext},\ \texttt{Quot.sound}\}$,
cleaner than the fixed base of the earlier sections) proves the bridge holds
verbatim for this strictly weaker demand, in three forms:

\begin{definition}[Continuation oracle; \texttt{CorrectContinuation}]
\label{def:contcont}
$\mathrm{CorrectContinuation}(f, T)$ iff for every requirement $(w,c)$ of
$T$ there exists a position $x$ with $f(w{+}{+}[c]) = \mathrm{some}(x)$,
$1 \le x \le |T|$, and $w{+}{+}[c]$ covered at $x$. The function $f$ is
required to do nothing on any other word.
\end{definition}

\begin{theorem}[Locate-one subsumption; \texttt{correctLocateOne\_imp\_cont}]
\label{thm:locimpcont}
If $f$ answers locate-one correctly on $T$ then $\mathrm{CorrectContinuation}(f,T)$.
\end{theorem}

Since requirement words occur, a correct locate-one oracle answers each fork
query with an occurrence of the continued word---which is exactly a covering
position. The continuation oracle is therefore the strictly weaker demand,
and every result below covers the locate-one case while standing on its own.

\begin{theorem}[Continuation bridge; \texttt{emitted\_suffixient\_of\_cont},
\texttt{chi\_le\_of\_cont\_oracle}]\label{thm:contbridge}
If $\mathrm{CorrectContinuation}(f,T)$ then the positions $f$ emits on
$T$'s requirement words form a suffixient set, and consequently
$\X(T) \le |\mathrm{emitted}(f,T)|$.
\end{theorem}

\begin{theorem}[Distinct-consultation; \texttt{chi\_le\_of\_cont\_oracle\_distinct}]
\label{thm:contdistinct}
If $\mathrm{CorrectContinuation}(f,T)$ then
$\X(T) \le \big|\mathrm{dedup}(\mathrm{emitted}(f,T))\big|$: across all
fork queries, $f$ must \emph{consult} at least $\X(T)$ \emph{distinct}
positions.
\end{theorem}

The distinct form is the honest statement of ``must retain'': an oracle that
re-emits one position $\X$ times has not consulted $\X$ positions. It
holds because duplicate removal preserves both suffixience and positional
membership (\texttt{suffixient\_dedup}, \texttt{dedup\_length\_le}), so the
deduplicated emitted set is itself a suffixient set of genuine text positions.

\begin{theorem}[Fixed-decoder index form; \texttt{chi\_le\_of\_realized\_cont}]
\label{thm:realizedcont}
For any decoder $\mathrm{dec}$ and index $D$: if the decoded answer function
$\mathrm{dec}(D)$ is a correct continuation oracle on $T$, then
$\X(T) \le \big|\mathrm{dedup}\big(\mathrm{emitted}(\mathrm{dec}(D),T)\big)\big|$.
\end{theorem}

This is ``any continuation-answering system must consult $\ge \X$
distinct positions'' in its precise, kernel-checked form. The
bridge-to-counting structure is now identical to the locate-one case
(Theorem~\ref{thm:bridge} $\to$ Theorem~\ref{thm:pigeon} $\to$
Theorem~\ref{thm:floor}); the continuation-counting analog of
Theorem~\ref{thm:floor}---replacing the locate-one demand with
$\mathrm{CorrectContinuation}$ in the family floor---is a mechanical
follow-on, since the incompatibility machinery of the witness-perturbation
family needs only that correct answers cover the requirements, which
continuation answers do by definition. We do not claim it as proven here.

\begin{theorem}[Pigeonhole; \texttt{family\_counting}]\label{thm:pigeon}
For a fixed decoder $\mathrm{dec}$, a duplicate-free family $F$ of texts
whose members pairwise admit no common correct answer function satisfies
$|F| \le 2^{s+1}-1$ for any $s$-bit index space serving all of them.
\end{theorem}

\paragraph{The witness-perturbation family.} For $k$ blocks of period $L$
and perturbation offsets $p_i \in [0,L/2]$, let
$T(p) = \mathrm{famText}(k,L,p)$ (block $i$ carries its marker at offset
$p_i$; the full definition is \texttt{LowerBound.lean}, which also proves
$2k \le \X(T) \le 3k-1$ (\texttt{chi\_fam\_bounds}) and
$n = k(L+1)$). Members with distinct offsets have pairwise incompatible
correct answers (\texttt{fam\_forced\_incompat}), and the half grid
$\{p : 2p_i \le L\}^k$ injects into texts (\texttt{famText\_inj\_grid}).

\begin{theorem}[Family floor; \texttt{fam\_floor\_chi}]\label{thm:floor}
For all $k\ge1$, $L\ge1$, base offsets $p_0$ in the half grid, and every
decoder $\mathrm{dec}$: if for every half-grid $p$ there exists an index
$D$ with $\mathrm{space}(D) \le s$ answering locate-one correctly on
$\mathrm{famText}(k,L,p)$, then
\[
\frac{\X\big(\mathrm{famText}(k,L,p_0)\big)\cdot
\log_2\!\big(k(L{+}1)\,/\,\X(\mathrm{famText}(k,L,p_0))\big)}{3} \;\le\; s+1 .
\]
\end{theorem}

\paragraph{What this does and does not establish.} It is the first space
lower bound parametrized by a repetitiveness measure for a pattern-matching
class (locate-one), and its rate is the natural tier
$\X\log(n/\X)$. It is \emph{family-specific}: the incompatibility
machinery is proven for the constructed family, and the general rung---every
text, or a measure-theoretic statement---remains open, as does the placement
of $\X$ against $\delta\log n$ that governs the general statement's final
form (\S\ref{sec:open}). The index model is fixed-decoder with bit-space
measurement; removing the fixed-decoder restriction is open. Since the
continuation bridge (\S\ref{sec:contbridge}) has the same shape as the
locate-one bridge, the continuation-counting analog of
Theorem~\ref{thm:floor}---the same bound for indexes answering
continuation queries only---is a mechanical follow-on; we record it as such,
not as proven. One skeleton
statement of the original proof plan (\texttt{fam\_oracle\_witness}) was
found \emph{refutable as locked} (counterexample $k{=}1,L{=}2,s{=}1$); the
intended content is instead proven as \texttt{fam\_floor\_half\_grid}, and
Theorem~\ref{thm:floor} does not depend on the refuted skeleton. We record
this because statement-lock discipline requires the counterexample to be
visible, not buried.

\section{The System}\label{sec:system}

The deployed tool (\texttt{xsa build}) is one command:
archive or FASTA/FASTQ/raw text in; a single container (\texttt{.sxi}) out,
holding the run-length BWT with its character counts, both packed SA-sample
sets (run tails and run heads), the anchors, the corpus names, and
$\delta$-compressed $\X$. Publication is fail-closed: multi-string corpora
require either byte-equal endpoint gates against an independently built
control or an in-flight audit that verifies sampled witnesses directly
against the source archive; corrupted sources and drifted tool binaries are
rejected before any scratch is touched (a 75-artifact hash manifest pins
every executable).

\paragraph{No materialization.} The pipeline never writes the text to disk.
The parse stage reads the compressed archive directly: we contributed an
AGC reader mode to the prefix-free parser (patch against
\texttt{marco-oliva/pfp}, PR materials prepared), which serves the canonical
extraction---uppercase reversed contigs in archive order, separated and
terminated by the reserved separator byte---from in-process archive queries
with thread-parallel prefetch. Measured steady reader throughput is
$440$--$535$\,MB/s at 16--96 threads on HPRC windows (73--87\% of the
file-based control), and the full-corpus parse projects to $6.65$\,h.
Downstream stages consume only parse-space and $r$-space artifacts; the
audit fetches its 10{,}000 witness ranges from the archive.

\paragraph{Stage inventory.} Parse (native archive read) $\to$ L2 parse
$\to$ prefix-free BWT front end with an additive endpoint tap (a 27-line
vendored patch emitting the run-tail SA samples the merge loop already
computes; the head samples are untouched, byte-identical to upstream) $\to$
endpoint normalization and seam repair (with refusal policy) $\to$ parse-space
aggregate computation (two-level fingerprint LCE, O(1) resolve from head
samples, streamed aggregate sweep) $\to$ container write and validation.

\section{Results}\label{sec:results}

\begin{table}[t]
\centering\small
\caption{Corpora and headline measurements. The 466-haplotype number is a
placeholder pending the running build; the historical BCR-frame value for
that corpus is $2{,}249{,}968{,}075$. Sources: Appendix~\ref{app:facts}.}
\begin{tabular}{lrrrr}
\toprule
Corpus & $n$ (bp) & strings & $\rparam$ (runs) & $\X$ \\
\midrule
yeast (single string) & $3.34\times10^9$ & 1 & $100{,}904{,}881$ & $85{,}404{,}240$ \\
yeast235 (from AGC) & $3.34\times10^9$ & 9{,}901 & $100{,}905{,}045$ & $85{,}404{,}336$ \\
k10 human subset & $3.02\times10^{10}$ & 865 & $1{,}859{,}825{,}862$ & $1{,}627{,}067{,}257$ \\
HPRC 466 & $1.40\times10^{12}$ & \emph{to measure} & \emph{to measure} & \hprcchi \\
\bottomrule
\end{tabular}
\end{table}

\paragraph{Yeast, born from scratch.} The single-string yeast corpus indexes
from zero in 26 minutes at 8.8\,GB peak: parse 79\,s, front end 624\,s,
endpoints 25\,s, aggregates 833\,s at 3.7\,GB, sweep 15\,s at 6\,MB, write
and validate 32\,s. The container is 1.80\,GB with $\X=85{,}404{,}240$
embedded ($\delta$-ratio $0.1299$); the fresh tails and run table are
byte-identical to the independently built pilot artifacts, and the fresh
heads byte-identical to the $\varphi^{-1}$ extractor's output against the
walk baseline. $\X/\rparam = 0.846$.

\paragraph{Three builds, one answer.} The 9{,}901-contig yeast235 corpus was
published three times from three independent toolsets (the original
seam-repair toolchain, the native-archive toolchain, and a fresh-checkout
sealed build under the distribution manifest): all three agree on all five
core container members and on $\X = 85{,}404{,}336$.

\paragraph{The boundary-delta identity, measured.} The single-string and
separated forms of the same yeast content differ by $+96$ in $\X$
($85{,}404{,}336$ vs $85{,}404{,}240$, $k = 9{,}901$ separators): the
separator-adjacent witnesses are the entire difference, a first measured
instance of the boundary-delta identity. At 30\,Gbp human scale the
canonical-frame $\X = 1{,}627{,}067{,}257$ differs from the historical
BCR-frame value $1{,}627{,}063{,}183$ by $+4{,}074$
($\approx 4.7$ per string, $k=865$). We record honestly that the
\emph{same-frame witness attribution} of this delta is not proven; the
statement-level theorem (separator insertion perturbs $\X$ by $O(k)$ terms)
is stated and locked with its obstruction documented, and both measured
deltas are consistent with it.

\paragraph{k10 from zero, twice.} The 30\,Gbp corpus was built twice
independently (control and gated); the complete publications are
byte-identical. Per-stage medians: parse 827\,s (9.7\,GB), L2 parse 54\,s,
front end 7{,}910--8{,}211\,s (111\,GB), endpoints 1{,}556--1{,}731\,s
(125\,GB), aggregates 8{,}829--10{,}155\,s (66\,GB), sweep 201--213\,s,
write 416--490\,s, validate 129--138\,s; peak 134\,GB at 16 threads. The
seam repair touched 7 classes / 1{,}161 rows at $1.86\times10^9$ runs.

\begin{table}[t]
\centering\small
\caption{Measured archive-reader throughput (steady state, identical
checksums against the file-based control in every row). Source:
\texttt{THROUGHPUT.md}.}
\begin{tabular}{lrrrr}
\toprule
Corpus & $j$ & AGC MB/s & file MB/s & ratio \\
\midrule
yeast235 & 16 & 470.7 & 647.0 & 72.7\% \\
yeast235 & 48 & 509.0 & 625.6 & 81.4\% \\
yeast235 & 96 & 534.8 & 615.7 & 86.9\% \\
HPRC windows & 16 & 509.1 & 692.7 & 73.5\% \\
HPRC windows & 48 & 464.3 & 666.8 & 69.6\% \\
HPRC windows & 96 & 437.6 & 730.1 & 59.9\% \\
\bottomrule
\end{tabular}
\end{table}

\paragraph{Cost per run, honestly decomposed.} At yeast scale the container
stores $17.9$ bytes per BWT run as built: $\approx5.0$ for the run-length
BWT and counts, $\approx4.0$ for packed tails, $\approx8.0$ for heads
(stored raw rather than packed---identified, not yet reclaimed),
$\approx0.88$ for the $\delta$-coded witnesses. The representation floor
for a locate-capable run-length index is $\approx12.7$\,B/run (two log-$n$
sample words plus the run table). What the \emph{proven} theorem
underwrites is different and smaller: Theorem~\ref{thm:floor} implies
$O(\X\log(n/\X))$ \Bits\ of decision-structure storage is
\emph{unavoidable}---at yeast, $\approx18.8$\,MB with the proven constant
($0.19$\,B/run; $0.56$\,B/run at the natural tier without the constant
$1/3$)---which underwrites about $4\%$ of the engineering floor. The
witness member as built ($88.7$\,MB $=0.88$\,B/run) sits within
$1.6\times$ of the natural tier. The honest reading: the samples and run
table are representation cost (known, packable); the witnesses are
decision-structure cost (provably irreducible below
$\X\log(n/\X)$ bits, up to the constant we have proven).

\section{The Fork Skeleton and the Reader}\label{sec:lm}

The continuation bridge has a reader-facing reading that deserves its own
statement, because it is the reason we believe $\X$ measures something about
\emph{corpora} rather than only about indexes.

\paragraph{The decision structure of a corpus.} Combine the two bridge forms
with the forced/fork split. Every position of $T$ is either \emph{forced}---a
unique continuation, bookkeeping, compressible without loss of
answerability---or part of a \emph{fork}: a right-maximal context with $\ge 2$
continuations, witnessed by some requirement. Theorem~\ref{thm:contdistinct}
says any system that answers continuation queries \emph{with positions} must
consult $\ge\X$ distinct fork-adjacent positions; it may discard
\emph{everything else}. The corpus's decision structure is thus precisely
$\X$-shaped: $n - \X$ bytes of inevitability plus a sparse skeleton of
genuine choices, and the theorem prices the skeleton at
$\Omega(\X\log(n/\X))$ bits (Theorem~\ref{thm:floor}, family-specific form).
At 466-haplotype scale the skeleton is $0.16\%$ of the corpus: a system that
has retained the forks has retained what cannot be compressed away, and a
system that has not, cannot answer.

\paragraph{The memorization-externalization thesis, now warranted.} A
language model trained on a corpus is, among other things, a system that
answers continuation queries. The design thesis this project's companion
specification pursues---\emph{memorize nothing, externalize the rulebook into
the index, spend network capacity on the policy of choosing among forks}---now
has its formal warrant: by Theorem~\ref{thm:contdistinct}, the
index-retained positions are not an optimization but a \emph{necessity} for
any position-consulting continuation system; and by the forced/fork split, a
model may offload the $n - \X$ inevitable bytes entirely. The bridge is
direct evidence for the position-consulting regime; whether an implicit
(weight-space) continuation system must obey the same bound is a
statement-lock candidate we do \emph{not} claim. The fork-conditioned
architecture (hard-masked next-byte prediction at the forks, free passage at
forced positions, witness pointers as episodic memory) is specified and
compute-audited in the companion document; its evaluation is future work.

\section{Related Work}\label{sec:related}

Our survey (verified against fetched sources; the full annotated map with
arXiv links is \texttt{LOWER\_BOUND\_PLAN.md}) found:

\textbf{No space lower bounds for pattern matching.} The r-index
\cite{r-index} achieves $O(r)$ words with optimal-time count/locate---where
``optimal'' is time-optimal \emph{within} that space; no $\Omega(r)$ space
bound for count or locate exists in anything we found. MONI
\cite{moni} finds MEMs in $O(r)$ space; the measure paper
\cite{measure-paper} proves $\gamma \le \X \le 2\overline r$ and that
suffixient sets are attractors (giving $O(\X\log(n/\X))$-space text
storage), and proves minimality \emph{among suffixient sets}; the
bridge-and-counting construction of \S\ref{sec:floor} is, as far as the
survey found, the first lower bound transferring that minimality to answer
functions of arbitrary (fixed-decoder) indexes.

\textbf{Attractors and the collapse program.} String attractors
\cite{attractors} lower-bound every dictionary-compression measure and
carry a time-space access tradeoff; since $\gamma\le\X$, attractor floors
transfer only up to the $\X$-vs-$\gamma$ gap (on de Bruijn-like families
$\X$ can exceed $\gamma$ asymptotically, which is exactly where our
witness-perturbation counting must do its own work). The suffix-array
collapse theorem \cite{collapse} gives full SA functionality in
$\delta$-space, and SODA 2026 equivalences \cite{soda26} map the tradeoff
landscape; whether pattern matching (not just SA access) collapses to
$\delta$-space is open and governs our general rung (\S\ref{sec:open}).

\textbf{MEMs in compressed space.} Navarro computes MEMs on repetitive
collections in $O(\delta\log n)$ words \cite{navarro-mems}; this is the
upper bound that links any $\X$-floor to the open
$\X$-vs-$\delta\log n$ placement: if $\X = O(\delta\log n)$ always, the
natural-tier floor is consistent; if not, an $\Omega(\X)$ floor for MEMs is
refuted and the floor delimits precisely which operations it governs.

\textbf{Construction tools.} Prefix-free parsing \cite{pfp}; the r-index
BWT construction and its front end (r-pfbwt); the AGC archive format
\cite{agc} whose native range queries our parser now consumes; ropebwt3
\cite{ropebwt3}, whose MEM output format our query layer emits byte-compatibly
(verified against the reference implementation built from source on 23
fixtures).

\section{Open Problems}\label{sec:open}
\begin{enumerate}[leftmargin=*]
\item \textbf{RunEdgeDominate.} The capstone's single remaining hypothesis:
every run-edge position is dominated by an emitted one. Zero measured
counterexamples; blocked on the emission rule's exact semantics.
\item \textbf{The general floor.} Theorem~\ref{thm:floor} is
family-specific; the general rung requires either a measure-theoretic
incompatibility construction or the negative resolution via
$\X$-vs-$\delta\log n$. The continuation-counting analog (same bound for
continuation-only indexes; the bridge shapes are identical, \S\ref{sec:contbridge})
is the nearest mechanical step.
\item \textbf{The implicit form.} Whether a weight-space (non-position-consulting)
continuation system---e.g.\ a trained language model without retrieval---must
obey the $\Omega(\X)$ bound is not claimed; a statement-lock candidate is
recorded in the companion specification, unproven.
\item \textbf{Minimality of $\X$ itself} among repetitiveness measures for
this operation class, and the exact constant of
Theorem~\ref{thm:floor} (currently $1/3$).
\item \textbf{The seam reordering lemma} (in-class spliced output equals
cyclic order): statement-locked, byte-warranted, unproven.
\item \textbf{The boundary-delta identity}: $\X$ of the $k$-fold
separator insertion differs from $\X$ of the base text by $O(k)$;
obstruction documented (position-shift/coverage coupling in the class
characterization).
\end{enumerate}

\section*{Acknowledgments}
The corrections ledger of this project---twelve entries at the time of
writing, every one caught before it reached a claim---is the real
acknowledgments section; the method is the coauthor.

\bibliographystyle{plain}
\bibliography{refs}

\appendix
\section{Claim-to-Artifact Map}\label{app:facts}
Every numeric and formal claim above, with its machine source, is listed in
\texttt{paper/FACTS.md} in the repository. The Lean statements quoted are in
\texttt{lean/}; the measurement records in \texttt{bit6/sxi\_logs/} and the
ledger \texttt{RESEARCH.md}.

\end{document}
